\subsection{A new exact optimum at fifty vertices}
The preceding incoming report leaves exact optima at large lengths not
divisible by eight undetermined. The following theorem settles its first
new case investigated here at which the alternating family fails. This is
a new result relative to the 46 inspected incoming archives; no worldwide
priority claim is made.

Write $a_i\in\{\pm1\}$ for the sign on the edge $\{i,i+1\}$ and
$b_i\in\{\pm1\}$ for the sign on $\{i,i+2\}$, with indices modulo $50$.
The triangle signs and Hamilton-cycle holonomy are
\[
 \tau_i=a_i a_{i+1}b_i,\qquad \alpha=\prod_{i=0}^{49}a_i.
\]
Vertex switching preserves these data, and they determine a labeled
switching class: switch to $a_0=\cdots=a_{48}=1$, $a_{49}=\alpha$,
then recover $b_i=\tau_i a_i a_{i+1}$.

\begin{theorem}[The fifty-vertex optimum]
Let $\rho_{50}$ be the minimum spectral radius over all independent edge
signings of $C_{50}(1,2)$. Put
\[
 B=(-,-,+,-,+,+,-,+),\qquad W=B^6(-,+).
\]
Precisely the switching classes with $\alpha=+1$ and triangle word a
cyclic rotation of $W$ or $-W$ minimize the spectral radius. These are
$100$ distinct labeled switching classes. Moreover,
\[
  \boxed{2.811568906156<\rho_{50}<2.811568906157.}
\]
An exact algebraic specification is
$\rho_{50}=\sqrt{y_{\max}}$, where $y_{\max}$ is the largest real root
of the degree-$25$ polynomial $Q$ recorded in the coefficient table below.
Every other switching class has spectral radius at least
$7029/2500=2.8116$.
\end{theorem}

\begin{proof}
Set $T=7029/2500$. For a triangle word $w$ of length $k$, let $H(w)$
be the $(k+2)$-vertex path with all length-one signs positive and
length-two signs given by $w$. Retain a word exactly when both
$TI-H(w)$ and $TI+H(w)$ are positive definite. Starting from the empty
word and extending only retained prefixes exhausts all passing words,
because each shorter path is a principal submatrix.

The exact Schur update on the last two vertices is
\[
 (a,b,c)\longmapsto
 \left(c-\frac{b^2}{a},\,-1+\frac{bt}{a},\,T-\frac1a\right),
\]
for appended sign $t$. Start with $(T,-1,T)$; the second shifted matrix
uses $-t$. The updated block passes precisely when $a'>0$ and
$a'c'>(b')^2$ in both cases. All calculations use rational arithmetic.
At triangle length $46$ there are $324$ passing words. Append all four
remaining signs to each of them and retain a length-$50$ word only if
all its cyclic length-$46$ windows pass. There are exactly two resulting
rotation classes, represented by $W$ and $-W$.

The exhaustive local path classification is checked by two independent
algorithms in the package. The first uses the Schur recurrence; the
second constructs determinants by summing signed permutation products,
using only integer arithmetic and the band width of the actual path
matrices. They agree on every prefix count and every surviving word.
The cyclic completion and full-matrix tests are exact finite calculations
in the first verifier. The complete lists and checks are supplied as
executable certificates.

Any consecutive $48$ vertices in $C_{50}(1,2)$ induce such a path:
the two endpoints have cyclic distance three, so no extra edge joins
them. Its matrix norm is bounded by the whole symmetric-matrix norm.
Consequently every signing with radius $<T$ must have one of these two
triangle words, up to rotation. There remain just two holonomies for
each word. Exact fraction-free elimination of the full matrices
$7029I\pm2500A$ shows that $\alpha=+1$ passes both Sylvester tests,
whereas $\alpha=-1$ fails. The first $49$ leading minors are positive
in each failed test, and the final determinant is negative. Thus no
additional signing can have radius $<T$.

Negating every edge changes every triangle sign to its negative and,
because $50$ is even, preserves $\alpha$. It takes $A$ to $-A$, so the
two surviving rotation classes have equal spectral radius. Both are
therefore global minimizers. Each word has primitive period $50$ and
their rotation orbits are disjoint, giving exactly $100$ labeled
switching classes.

The characteristic polynomial of either minimizing representative is
$\det(xI-A)=Q(x^2)$. The supplied checker computes it from integer traces
$\operatorname{tr}(A^j)$ by Newton's identities. For
$r_-=2811568906156/10^{12}$ and
$r_+=2811568906157/10^{12}$, exact leading-minor sign counts show
that $r_-I-A$ and $r_-I+A$ each have exactly one negative eigenvalue,
while both shifts at $r_+$ have none. This proves the strict interval,
and also shows that the interval contains exactly one positive
eigenvalue. The characteristic polynomial specifies the value exactly.
\end{proof}

The alternating family's best value is
$2\sqrt{\cos^2(\pi/50)+\cos^2(2\pi/50)}>2.814$.
For an exact comparison, use $\pi<22/7$ and
$\cos u\ge1-u^2/2$: setting $v=11/175$, rational arithmetic gives
\[
 4\left((1-v^2/2)^2+(1-2v^2)^2\right)>(2814/1000)^2.
\]
Thus the improvement over alternation is strict and does not rely on
numerical eigenvalue estimates.

\paragraph{Polynomial coefficients.}
Write $Q(y)=\sum_{j=0}^{25}q_jy^j$. Its coefficients are
\begin{center}\small
\begin{tabular}{rr@{\qquad}rr}
\hline
$j$&$q_j$&$j$&$q_j$\\\hline
25&1&12&-42556190827574\\
24&-100&11&95424283730427\\
23&4702&10&-173149057931144\\
22&-138186&9&251857479144425\\
21&2846105&8&-290088695299758\\
20&-43660212&7&260508422825773\\
19&517548579&6&-178938476667148\\
18&-4854567112&5&91816438267915\\
17&36601189842&4&-34160308422738\\
16&-224095627758&3&8848896828974\\
15&1121189638159&2&-1499411878484\\
14&-4598060887914&1&148268409185\\
13&15463123646497&0&-6447640168\\\hline
\end{tabular}
\end{center}
No irreducibility assertion is needed. Exact finite enumeration and
the analytic compression argument establish the theorem; this result
has not been formalized in Lean.
